\section{1B Agent Model Architecture and Comparison Plan}

\begin{notesquote}
Version: v0.3

Status: the core formulas of Utility Consolidation Memory v1 are frozen; the full 1B backbone is not yet frozen.

Goal: train from scratch a small Agent model of about 1B parameters, targeting tool calling, simple code, basic mathematics, and long-task state retention.
\end{notesquote}

\subsection{Project Goals}

The model is positioned as a small Agent controller rather than a general chat model that covers all knowledge. Its core capabilities include:

\begin{enumerate}
  \item Understanding Chinese and English tasks;
  \item Generating structurally correct tool calls;
  \item Writing and repairing simple code;
  \item Performing basic mathematical reasoning;
  \item Interpreting tool results and performing limited error correction from failures;
  \item Maintaining goals, constraints, state updates, and key tool results in long tasks under a fixed state budget;
  \item Producing concise, verifiable final answers.
\end{enumerate}

Exact computation and program execution are delegated to external tools. The model focuses on learning task decomposition, tool selection, argument construction, result interpretation, stopping conditions, and memory trade-offs under limited capacity.

\subsubsection{Scope of the First Version}

\begin{itemize}
  \item Basic arithmetic, algebra, ratios, unit conversion, and simple word problems;
  \item Small Python programs and basic Bash/Shell tasks;
  \item Reading small code files and repairing code based on unit test failures;
  \item Calculator, Python, Shell, file, and test tool calls;
  \item Controlled Agent or state-tracking trajectories of 8--64 steps;
  \item Final state queries after variables, constraints, file state, and tool results have been overwritten multiple times;
  \item Selective recall of key events once they fall outside the local window.
\end{itemize}

\subsubsection{Out of Scope for the First Version}

\begin{itemize}
  \item Large repository-level software engineering;
  \item Difficult competition mathematics;
  \item Permanent personalized memory across users and across sessions;
  \item Multi-Agent collaboration, GUI automation, and multimodality;
  \item Lossless retention of histories of arbitrary length;
  \item Surpassing Full Attention across the board under short-context or unlimited-compute conditions.
\end{itemize}

\subsection{Current Architecture Status}

\subsubsection{Frozen: Utility Consolidation Memory v1}

Only one candidate enters long-term memory competition at a time:

\[
\mathcal R_t=M_t\cup\{c_t\}.
\]

During training, counterfactual future loss after the decision generates the utility labels:

\[
u_{t,i} =
\mathcal L_{\bar\theta,F}
\left(Y_t\mid B_t,\operatorname{mask}_i(\mathcal R_t)\right) -
\mathcal L_{\bar\theta,F}
\left(Y_t\mid B_t,\mathcal R_t\right).
\]

The predictor may only use information already observed at decision time:

\[
\hat u_{t,i}=f_\phi(m_i,g_t,B_t,\mathcal R_t).
\]

When capacity is exceeded, only the single record with the lowest predicted utility is evicted:

\[
M_{t+1} =
\mathcal R_t\setminus
\left\{ \arg\min_i\hat u_{t,i}\right\}.
\]

See \href{utility-consolidation-memory-v1.md}{Utility Consolidation Memory v1} for the complete definition.

\subsubsection{Explicitly Excluded from v1}

\begin{itemize}
  \item Explicit age, occupancy, or time decay;
  \item Novelty, surprise, or conflict gates;
  \item Event association graphs;
  \item soft write;
  \item Interpolation, rewriting, or merging of memory content;
  \item Gumbel-Softmax or differentiable Top-K;
  \item Backpropagation or future information access at inference time.
\end{itemize}

\subsubsection{Still Not Frozen}

\begin{lstlisting}[style=notes]
Parameter scale: about 1.0B--1.3B
Model type: autoregressive text generation model
Training precision: BF16
Primary languages: Chinese and English
Primary domains: tool calling, simple code, basic mathematics, long-task state retention
Deployment target: inference on a single 16GB--48GB GPU
\end{lstlisting}

The following still require controlled experiments:

\begin{itemize}
  \item Whether the backbone uses Full/Local Attention, linear attention, or a hybrid structure;
  \item The number of layers, hidden size, FFN width, and the layer position for memory reads;
  \item The local window $W$;
  \item The long-term capacity $S$, the waiting period $D$, and the supervision horizon $F$;
  \item Event record granularity, encoding dimension, and how the original text is retained;
  \item The concrete implementation of the Utility Predictor;
  \item Positional encoding, Tokenizer, data mixture, and training tokens.
\end{itemize}

Freezing the memory formulas does not mean that the final 1B backbone has been decided.

\subsection{Public Models and Architecture References}

\subsubsection{TinyLlama-1.1B: A Standard Small Transformer}

\begin{table}[htbp]\centering
\begin{tabular}{lr}\toprule
Item & Configuration\\
\midrule
Parameter count & about 1.1B\\
Layers & 22\\
Hidden size & 2048\\
FFN size & 5632\\
Query/KV heads & 32/4\\
Attention & GQA\\
Activation/normalization & SwiGLU/RMSNorm\\
Positional encoding & RoPE\\
Vocabulary/context & 32K/2048\\
\bottomrule
\end{tabular}
\end{table}

References: \href{https://github.com/jzhang38/tinyllama}{TinyLlama official repository}, \href{https://github.com/jzhang38/TinyLlama/blob/main/PRETRAIN.md}{pretraining notes}.

\subsubsection{OLMo-1B: An Open Training Reference}

Used to compare training curves, data efficiency, and language model capability at a similar number of tokens. References: \href{https://github.com/allenai/OLMo}{OLMo official repository}, \href{https://huggingface.co/allenai/OLMo-1B}{OLMo-1B model card}.

\subsubsection{RWKV-7: A Fixed Recurrent State Reference}

RWKV-7 maintains a fixed recurrent state with a generalized Delta Rule and vector gating and serves as a continuous-state control. It is used to test whether discrete, addressable event records really outperform a state matrix that continuously compresses history.

References: \href{https://arxiv.org/abs/2503.14456}{RWKV-7}, \href{https://github.com/BlinkDL/RWKV-LM}{official repository}.

\subsubsection{Kimi Linear / KDA: A Linear Attention Reference}

The core state update of KDA can be summarized as:

\[
S_t=(D_t-a_tb_t^\top)S_{t-1}+k_tv_t^\top.
\]

It is used to compare the quality and efficiency differences between continuous matrix state, hybrid Full Attention, and discrete event retention. Reference: \href{https://arxiv.org/abs/2510.26692}{Kimi Linear}.

Kimi K3 is a large-scale combination of KDA, Gated MLA, Attention Residuals, and Stable LatentMoE and does not serve as a strict 1B capability baseline.

\subsubsection{Other Capability and Hybrid Architecture References}

\begin{itemize}
  \item \href{https://arxiv.org/abs/2411.13676}{Hymba-1.5B}: small-model hybrid architecture, caching, and long-context efficiency;
  \item \href{https://arxiv.org/abs/2409.12186}{Qwen2.5-Coder-1.5B}: simple code and Code Agent capability;
  \item \href{https://huggingface.co/HuggingFaceTB/SmolLM2-1.7B/blob/main/README.md}{SmolLM2-1.7B}: small-model data engineering, mathematics, instruction following, and Function calling.
\end{itemize}

\subsubsection{Adjacent Memory Methods}

\begin{itemize}
  \item \href{https://arxiv.org/abs/2506.15841}{MEM1}: a joint reasoning--memory state of fixed capacity;
  \item \href{https://arxiv.org/abs/2605.21768}{Memory-R2}: local and global credit assignment for long-range memory operations;
  \item \href{https://arxiv.org/abs/2501.00663}{Titans}: Attention and test-time neural long-term memory;
  \item \href{https://arxiv.org/abs/2404.07143}{Infini-attention}: local attention and bounded compressed memory;
  \item \href{https://arxiv.org/abs/2207.06881}{Recurrent Memory Transformer}: recurrent memory tokens across segments.
\end{itemize}

\subsection{Comparison Responsibilities}

\begin{table}[htbp]\centering
\begin{tabular}{llr}\toprule
Subject & Main comparison content & Strict architectural control\\
\midrule
TinyLlama-1.1B & Standard small GQA, parameter efficiency & No\\
OLMo-1B & Open training process, similar token budget & No\\
RWKV-7 & Fixed continuous state, throughput, state tracking & Yes in the controlled implementation\\
KDA & Linear attention, continuous matrix state & Yes in the controlled implementation\\
Hymba/Qwen/SmolLM2 & Hybrid, code, mathematics, and Agent capability & No\\
FIFO/random eviction & Lower bound of memory policies & Yes\\
Future Oracle eviction & Offline upper bound of the Reader and the record representation & Yes\\
Predicted utility eviction & Deployable policy of this project & Yes\\
\bottomrule
\end{tabular}
\end{table}

\subsection{Core Research Hypothesis}

This project does not claim to surpass Transformers or linear attention across the board. The claim to be validated is:

\begin{notesquote}
Under the same resident state budget, for trajectories where key events are sparse, tasks are very long, and state is overwritten multiple times, a single eviction policy trained with post-decision counterfactual future loss retains future-useful information more reliably than time-based eviction and continuously compressed state.
\end{notesquote}

\subsubsection{Relative to Full/Local Transformers}

The KV state of a full-context Transformer grows with the total length:

\[
O(Td).
\]

This project uses a local window $W$ and a long-term capacity $S$:

\[
O((W+S)d).
\]

The difference from a sliding window is not whether the state is bounded, but that old records are evicted by predicted future value rather than by pure temporal order.

\subsubsection{Relative to Linear or Recurrent State Models}

The advantages of linear attention and RWKV/KDA are a fixed state and efficient per-token updates. The hypothesized advantages of this project are:

\begin{itemize}
  \item Event records keep independent identities;
  \item A single record can be addressed, masked, and deleted;
  \item History does not have to be continuously mixed into the same matrix state;
  \item Memory trade-offs have direct counterfactual future utility supervision;
  \item Delayed decisions can use evidence already observed after an event occurs.
\end{itemize}

This project may be at a disadvantage in the pure throughput of linear models, in dense sequence modeling, and on tasks whose future usage is highly unpredictable.

\subsection{Fair Comparison Rules}

\subsubsection{Public Side-by-Side Comparison and Controlled Experiments Are Kept Separate}

Public models are used to answer final capability level and deployment efficiency; controlled small models are used for architectural attribution. Public models with different Tokenizers, data, and training tokens cannot be used to prove that utility memory is effective.

\subsubsection{Parameters and State Budget}

Controlled models match total parameters, non-Embedding parameters, or training FLOPs, and report the other items in full. The state budget must include:

\[
\text{local KV}
+\text{short-term buffer}
+\text{long-term records}
+\text{metadata}.
\]

If the short-term buffer is fully contained in the local window, this should be stated explicitly to avoid double counting.

\subsubsection{Cross-Tokenizer Metrics}

Perplexity is not compared directly; Bits Per Byte is used uniformly:

\[
\operatorname{BPB} =
\frac{\operatorname{NLL}}
{N_{\mathrm{bytes}}\ln2}.
\]

Report Chinese, English, code, mathematics, and API/JSON BPB separately.

\subsection{Evaluation Metrics}

\subsubsection{General, Mathematics, Code, and Agent}

\begin{itemize}
  \item Chinese and English held-out BPB;
  \item Accuracy on basic mathematics and tool-assisted mathematics;
  \item Python pass@1, syntax correctness rate, and unit test pass rate;
  \item Success rate of repairing code from error logs;
  \item Tool-call JSON validity rate;
  \item Tool selection and argument construction accuracy;
  \item Multi-step task success rate, failure recovery rate, and final answer correctness rate.
\end{itemize}

\subsubsection{Long-Range Memory}

\begin{itemize}
  \item Fact recall across the local window;
  \item Latest-value correctness rate after multiple state overwrites;
  \item Stale-value misuse rate;
  \item Performance degradation as irrelevant tool output increases;
  \item Recall under different dependency distances $D,F$;
  \item Record eviction accuracy;
  \item Rank correlation between the Utility Predictor and the Oracle;
  \item Selection regret of the predicted policy relative to the Oracle;
  \item Transfer of the memory policy on a fully held-out domain.
\end{itemize}

\subsubsection{Efficiency}

\begin{itemize}
  \item Prefill and Decode tokens/s;
  \item First-token, single-token, and consolidation latency;
  \item Peak GPU memory and total state size;
  \item Number of consolidations per thousand tokens;
  \item FLOPs for counterfactual label generation;
  \item Total training tokens/s, per-step time, and MFU.
\end{itemize}

Offline counterfactual label cost must not be counted toward online latency at inference time; training cost must be reported in full.

\subsection{Experiment Stages}

\subsubsection{Stage 0: Selector Validation Without a Language Model}

Validate on synthetic state machines and key-value overwrite tasks:

\begin{enumerate}
  \item That the single eviction formula and the mask implementation are correct;
  \item Whether Future Oracle significantly exceeds FIFO;
  \item Whether the single eviction explanation holds when records are redundant or complementary;
  \item The effect of candidate order on the greedy policy.
\end{enumerate}

If Oracle does not exceed FIFO, do not proceed to larger-model experiments.

\subsubsection{Stage 1: 20M--50M Reader and Utility Predictor}

\begin{enumerate}
  \item First train the Reader to use external records;
  \item Freeze or EMA the Reader/event encoder;
  \item Generate a small number of exact counterfactual labels;
  \item Train the Utility Predictor;
  \item Compare FIFO, random, predicted utility, and Oracle.
\end{enumerate}

\subsubsection{Stage 2: 100M--150M Controlled Architecture Experiments}

Match Tokenizer, data order, parameters/FLOPs, batch, optimizer, training tokens, and random seeds. Compare:

\begin{itemize}
  \item Local Transformer + FIFO;
  \item Local Transformer + Utility Consolidation Memory;
  \item One RWKV/KDA-style continuous-state baseline;
  \item A Full Attention upper bound when necessary.
\end{itemize}

\subsubsection{Stage 3: 300M--400M Scaling Validation}

Scale only the approaches that simultaneously satisfy the following conditions:

\begin{itemize}
  \item Oracle has a clear upper bound;
  \item The predicted policy is significantly close to Oracle;
  \item It exceeds FIFO and the continuous-state baseline at the same state budget;
  \item Label cost is still acceptable.
\end{itemize}

\subsubsection{Stage 4: Final Model of About 1B Parameters}

In the end only one candidate of about 1B parameters is trained. The Transformer, linear, or hybrid backbone is determined by the quality--efficiency curves of stages 2--3.

\subsection{Draft Agent Tool Protocol}

The first version is expected to support:

\begin{lstlisting}[style=notes]
calculator
python
shell
read_file
apply_patch
run_tests
\end{lstlisting}

Initial limits:

\begin{lstlisting}[style=notes]
max_agent_steps: 64
max_retries_per_step: 2
python_timeout_seconds: 10
shell_timeout_seconds: 10
max_tool_output_chars: 12000
\end{lstlisting}

Short capability evaluations keep an 8-step subset; memory evaluations extend progressively to 16, 32, and 64 steps.

\subsection{Results Recording Template}

\subsubsection{Memory Selection}

\begin{table}[htbp]\centering
\begin{tabular}{lrrrrr}\toprule
Policy & Total state budget & Final success rate & Stale-value misuse rate & Selection Regret & Consolidation latency\\
\midrule
FIFO & & & & & \\
Random & & & & & \\
Predicted Utility & & & & & \\
Future Oracle & & & & & \\
\bottomrule
\end{tabular}
\end{table}

\subsubsection{Architecture Quality and Efficiency}

\begin{table}[htbp]\centering
\begin{tabular}{lrrrrrr}\toprule
Model & Parameters & Training Tokens & BPB & Agent success rate & Decode tok/s & State size\\
\midrule
Local Transformer & & & & & & \\
Linear/Recurrent & & & & & & \\
Utility Memory & & & & & & \\
Full Attention Upper Bound & & & & & & \\
\bottomrule
\end{tabular}
\end{table}

\subsection{Current Decisions and Stopping Conditions}

Decided:

\begin{itemize}
  \item Positioning as a small Agent model of about 1B parameters;
  \item Tool calling, simple code, basic mathematics, and long-task state retention;
  \item Not training a full conventional 1B baseline again;
  \item Public models carry the final capability reference;
  \item Small-scale controlled experiments carry the architectural attribution;
  \item The core formulas of Utility Consolidation Memory v1;
  \item Single-candidate admission, single-record eviction, and immutable event records;
  \item The first version adds no other memory gates or association modules.
\end{itemize}

Still undecided:

\begin{itemize}
  \item The full backbone sequence-mixing structure;
  \item Depth, width, and the layer position for memory reads;
  \item $W,S,D,F$;
  \item Event record encoding;
  \item Label sampling rate and the teacher update schedule;
  \item Tokenizer, data mixture, and final training tokens.
\end{itemize}

Stopping conditions:

\begin{enumerate}
  \item Future Oracle does not exceed FIFO at the same budget;
  \item Predicted utility cannot approach Oracle significantly over the long run;
  \item Under the same state budget, it does not exceed the sliding-window or continuous-state baseline;
  \item Label cost cannot be reduced to an acceptable range through sampling, caching, or distillation.
\end{enumerate}
